% drafts/appendix-a-exactness.tex — Proof of Algebraic Exactness

We prove \cref{thm:exactness}: for Hamiltonian coefficients
$c_k \in \Qi$, the encoded Pauli-string coefficients $d_m$ satisfy
$d_m \in \Qi$, and no floating-point rounding is introduced by phase
tracking.

\paragraph{Step 1: Ladder $\to$ Majorana.}
Each ladder operator is decomposed as
\begin{align}
  \ad{j} &= \tfrac{1}{2}(c_j - i\,d_j), &
  \an{j} &= \tfrac{1}{2}(c_j + i\,d_j),
\end{align}
where $c_j$ and $d_j$ are Majorana operators.  The coefficients are
$\tfrac{1}{2}$ and $\pm\tfrac{i}{2}$, both elements of
$\tfrac{1}{2}\Qi$.

\paragraph{Step 2: Majorana $\to$ Pauli string.}
Each Majorana operator $c_j$ or $d_j$ is a single Pauli string
with coefficient $\pm 1$ or $\pm i$ (an element of $\Zi$).  This
is a property of the encoding map, which assigns Pauli labels
$X, Y, Z$ along a tree path and multiplies only elements of $\Zi$.

\paragraph{Step 3: Pauli string multiplication.}
Two $n$-qubit Pauli strings are multiplied qubit-by-qubit.  At each
position, the product $\sigma_a \cdot \sigma_b = \phi_{ab} \,
\sigma_{a \oplus b}$ is computed via the finite table
(\cref{eq:pauli-product}).  The phase $\phi_{ab} \in \Zi$ is
accumulated via the group law of $\Zi$.

At the end of the $n$-qubit product, the accumulated phase
$\phi = \prod_{k=1}^{n} \phi_k$ is an element of $\Zi$.  It is
folded into the coefficient via:
\begin{itemize}
  \item $\phi = +1$: no change.
  \item $\phi = -1$: negate the coefficient ($c \mapsto -c$).
  \item $\phi = +i$: apply $c = a + bi \mapsto -b + ai$.
  \item $\phi = -i$: apply $c = a + bi \mapsto b - ai$.
\end{itemize}
Each of these operations is exact in IEEE~754 for any finite
\texttt{double} inputs: negation and component swapping introduce
no rounding.

\paragraph{Step 4: Coefficient arithmetic.}
When the input coefficients $c_k$ are elements of $\Qi$, they have
the form $c_k = \frac{p + qi}{r}$ where $p, q, r \in \mathbb{Z}$.
In practice, all encoding coefficients are multiples of
$\frac{1}{2}$, so we work in $\frac{1}{2}\mathbb{Z}[i]$.

Multiplication of $\frac{1}{2}\mathbb{Z}[i]$ elements produces
$\frac{1}{4}\mathbb{Z}[i]$ elements; after $k$ multiplications,
coefficients lie in $\frac{1}{2^k}\mathbb{Z}[i]$.  These are all
exactly representable as \texttt{double} pairs for the system sizes
of interest ($k \le 50$, giving denominators $\le 2^{50}$ which are
exact in 53-bit mantissa).

\paragraph{Step 5: Like-term combination.}
Terms with the same Pauli signature are combined by adding their
coefficients: $d_m = \sum_{k \in S_m} d_k$.  This is the only
place where floating-point addition occurs.  For $\Qi$ inputs,
we are adding elements of $\frac{1}{2^k}\mathbb{Z}[i]$, which is
exact when the sum does not exceed $2^{53}$ in magnitude (always
true for physically meaningful Hamiltonians).

Conversely, the signature matching is performed by string comparison,
which is exact.  Two Pauli strings are combined if and only if they
are identical operator-by-operator.  No approximate comparison is
used.

\paragraph{Conclusion.}
The composition of Steps 1--5 maps $\Qi$ coefficients to $\Qi$
coefficients with no rounding in phase tracking.  The only potential
source of rounding is Step~5 (coefficient addition), which is exact
for the denominator ranges arising in encoding practice.
\qed
